\documentclass[../main]{subfiles}
\ifSubfilesClassLoaded{\setcounter{chapter}{6}}{}
\begin{document}

\chapter{Physics, Collision, dan Rigidbody}

\begin{subcpmk}
  \item Sub-CPMK 5.1: Mengkonfigurasi Rigidbody dan Collider, mengimplementasikan OnCollisionEnter/OnTriggerEnter
  \item Sub-CPMK 5.2: Membedakan trigger dan collision dalam implementasi game
\end{subcpmk}

\section{Camera Follow}

Setelah player bergerak (Bab VI), kamera perlu mengikuti. Bab ini membahas pola camera follow. Kamera perlu mengikuti posisi bola agar pemain selalu melihat dari belakang-atas. Buat script \class{CameraController}:

\begin{unitycode}
using UnityEngine;

public class CameraController : MonoBehaviour
{
    public GameObject player;
    private Vector3 offset;

    void Start()
    {
        offset = transform.position - player.transform.position;
    }

    void LateUpdate()
    {
        transform.position = player.transform.position + offset;
    }
}
\end{unitycode}

\class{LateUpdate} dipanggil setelah semua Update. Dengan demikian, kamera bergerak setelah player bergerak di frame yang sama---menghindari jitter (getaran visual). Variabel \texttt{offset} menyimpan jarak awal antara kamera dan player; offset tersebut dijaga konstan agar sudut pandang tetap sama.


\section{Collision Detection dan Respon}

\begin{learningoutcomes}
\item Mengimplementasikan OnCollisionEnter dan OnTriggerEnter
\item Menghancurkan objek saat collision
\item Menerapkan damage system
\end{learningoutcomes}

\subsection{Event Collision}

\begin{lstlisting}[language={[Sharp]C}]
void OnTriggerEnter2D(Collider2D other)
{
    if (other.CompareTag("Enemy"))
    {
        Destroy(gameObject);
    }
}
\end{lstlisting}

\subsection{Implementasi Space Survivor}

Bullet-OnTriggerEnter: Destroy bullet dan asteroid, tambah score. Player-OnCollisionEnter2D: Game Over. Tag GameObject untuk identifikasi (Enemy, Bullet, Player).


\begin{aktivitas}
  \item Implementasikan collision Bullet-Asteroid (Destroy, AddScore)
  \item Implementasikan Player-Asteroid collision (Game Over)
  \item Gunakan Tag untuk identifikasi GameObject
\end{aktivitas}

\begin{latihan}
  \item Jelaskan syarat terjadinya collision!
  \item Apa perbedaan OnTriggerEnter dan OnCollisionEnter?
  \item Kapan menggunakan Rigidbody Kinematic?
\end{latihan}

\begin{asesmen}
\textbf{Praktik}: Implementasi collision detection lengkap untuk Space Survivor.
\textbf{Rubrik}: Lihat Lampiran A
\end{asesmen}

\begin{checklist}
  \item Saya memahami Rigidbody dan Collider
  \item Saya dapat mengimplementasikan collision detection
  \item Saya dapat membedakan Trigger dan Collision
\end{checklist}

\begin{rangkuman}
Rigidbody untuk fisika, Collider untuk deteksi. OnTriggerEnter untuk trigger (peluru), OnCollisionEnter untuk fisik. Gunakan Tag untuk identifikasi.
\end{rangkuman}

\ifSubfilesClassLoaded{
  \renewcommand{\bibname}{Daftar Pustaka}
  \bibliographystyle{plain}
  \bibliography{../references}
}{}
\end{document}
